\documentclass[a4paper]{arrowhead}

\usepackage[yyyymmdd]{datetime}
\usepackage{etoolbox}
\usepackage[utf8]{inputenc}
\usepackage[T1]{fontenc}
\usepackage{multirow}
\usepackage{longtable}
\usepackage{tikz}
\usetikzlibrary{positioning,shapes.multipart,arrows.meta}

\renewcommand{\dateseparator}{-}
\setlength{\emergencystretch}{3em}
\lstset{language=HTTP,breakatwhitespace=false,columns=fullflexible}

%% Special references
\newcommand{\mref}[2][]{{\textcolor{ArrowheadPurple}{\hyperref[sec:model:#2]{\ifstrempty{#1}{#2}{#1}}}}}
\newcommand{\oref}[1]{{\textcolor{ArrowheadBlue}{\hyperref[sec:operations:#1]{#1}}}}
\newcommand{\pdef}[1]{{\textcolor{ArrowheadGrey}{#1\label{sec:model:primitives:#1}\label{sec:model:primitives:#1s}\label{sec:model:primitives:#1es}}}}
\newcommand{\pref}[1]{{\textcolor{ArrowheadGrey}{\hyperref[sec:model:primitives:#1]{#1}}}}

% as template_sd.tex; the input and output arguments are passed already formatted (\mref or \pref)
\newrobustcmd\osubsection[3]{
  \addtocounter{subsection}{1}
  \addcontentsline{toc}{subsection}{\protect\numberline{\thesubsection}operation \textcolor{ArrowheadBlue}{#1}}
  \subsection*{
    \thesubsection\quad
    operation
    \textcolor{ArrowheadBlue}{#1}
    (\notblank{#2}{#2}{})
    \notblank{#3}{: #3}{}
  }
  \label{sec:operations:#1}
}
\newrobustcmd\msubsection[2]{
  \addtocounter{subsection}{1}
  \addcontentsline{toc}{subsection}{\protect\numberline{\thesubsection}#1 \textcolor{ArrowheadPurple}{#2}}
  \subsection*{\thesubsection\quad#1 \textcolor{ArrowheadPurple}{#2}}
  \label{sec:model:#2} \label{sec:model:#2s} \label{sec:model:#2es}
}
%%

\begin{document}
\sloppy

%% Arrowhead Document Properties
\ArrowheadTitle{blacklistManagement}
\ArrowheadServiceID{blacklistManagement}
\ArrowheadType{Service Description}
\ArrowheadTypeShort{SD}
\ArrowheadVersion{5.2.2}
\ArrowheadDate{2026-10-30}
\ArrowheadAuthor{Jerker Delsing}
\ArrowheadStatus{RELEASE CANDIDATE}
\ArrowheadContact{jerker.delsing@ltu.se}
\ArrowheadFooter{\href{www.arrowhead.eu}{www.arrowhead.eu}}
\ArrowheadSetup
%%

%% Front Page
\begin{center}
  \vspace*{1cm}
  \huge{\arrowtitle}

  \vspace*{0.2cm}
  \LARGE{\arrowtype}
  \vspace*{1cm}

  \Large{Service ID: \textit{"\arrowid"}}

  \vspace*{\fill}

  % Front Page Image

  \vspace*{1cm}
  \vspace*{\fill}

  % Front Page Abstract
  \begin{abstract}
    This document describes the \texttt{blacklist\allowbreak{}Management} service, version 1.0.0, produced by the Blacklist support system of Arrowhead 5.2.2.
  \end{abstract}

  \vspace*{1cm}
\end{center}

\newpage
%%

%% Table of Contents
\tableofcontents
\newpage
%%

\section{Overview}
\label{sec:overview}

This document describes the \texttt{blacklist\allowbreak{}Management} service, which provides operator management of the blacklist: adding systems with a reason and expiry, querying entries and inactivating them.
It is produced by the Blacklist system \cite{bl-sysd} in service version 1.0.0.

The rest of this document is organized as follows.
In Section \ref{sec:operations}, we describe the abstract message operations provided by the service.
In Section \ref{sec:model}, we end the document by presenting the data types used by the mentioned operations.
The concrete interfaces are specified in the HTTP/TLS/JSON IDD \cite{bl-blacklistManagement-http-idd} and the MQTT/TLS/JSON IDD \cite{bl-blacklistManagement-mqtt-idd}.

\subsection{Significant Prior Art}

New in Arrowhead 5.

\subsection{How This Service Is Meant to Be Used}

Operators, the onboarding procedure (off-boarded systems are blacklisted) and security tooling add entries with \texttt{create}, inspect them with \texttt{query} and lift them with \texttt{remove}.

\subsection{Important Delimitations and Dependencies}
\label{sec:delimitations}

Entries are inactivated, not deleted, so the history is kept. Access follows the management policy of the Blacklist.

\newpage

\section{Service Interface}
\label{sec:operations}

This section describes the interfaces to the \texttt{blacklist\allowbreak{}Management} service.
In particular, each subsection names an abstract operation, an input type and an output type, in that order.
The input type is named inside parentheses, while the output type is preceded by a colon.
Input and output types are only denoted when accepted or returned, respectively, by the interface in question.

All abstract data types named in this section are defined in Section \ref{sec:model}.

\begin{figure}[ht!]
  \centering
  \begin{tikzpicture}[font=\small]
    \node[draw=ArrowheadBlue, very thick, rectangle split, rectangle split parts=2, align=center,
          rectangle split part fill={ArrowheadBlue!12,white}, text width=9cm, inner sep=6pt] (b)
      {\guillemotleft{}Service\guillemotright{}\\\textbf{blacklistManagement}
        \nodepart{second} \texttt{query}\\ \texttt{create}\\ \texttt{remove}};
  \end{tikzpicture}
  \caption{SysML-style block of the \texttt{blacklist\allowbreak{}Management} service and its operations.}
  \label{fig:SD}
\end{figure}

The following interface operations are available.

\osubsection{query}{\mref{BlacklistQueryRequest}}{\mref[Blacklist\allowbreak{}Entry\allowbreak{}List\allowbreak{}Response]{BlacklistEntryListResponse}}

Returns a page of entries matching the filter.

\osubsection{create}{\mref[Blacklist\allowbreak{}Create\allowbreak{}List\allowbreak{}Request]{BlacklistCreateListRequest}}{\mref[Blacklist\allowbreak{}Entry\allowbreak{}List\allowbreak{}Response]{BlacklistEntryListResponse}}

Adds blacklist entries with reason and optional expiry.

\osubsection{remove}{\pref{List}$<$\allowbreak{}\pref{Name}$>$}{}

Inactivates the entries of the named systems.

\newpage

\section{Information Model}
\label{sec:model}

Here, all data objects that can be part of the \texttt{blacklist\allowbreak{}Management} service are listed in alphabetic order.
Note that each subsection, which describes one type of object, begins with the \textit{struct} keyword, which is used to denote a collection of named fields, each with its own data type.
As a complement to the explicitly defined types in this section, there is also a list of implicit primitive types in Section \ref{sec:model:primitives}, which are used to represent things like hashes and identifiers.

\msubsection{struct}{BlacklistCreateListRequest}

Entries to add to the blacklist.


\noindent\begin{tabularx}{\textwidth}{| >{\raggedright\arraybackslash}p{3.9cm} | >{\raggedright\arraybackslash\small}p{5.3cm} | X |} \hline
\rowcolor{gray!33} Field & Type & Description \\ \hline
entities & \pref{List}$<$\allowbreak{}\mref{BlacklistCreateRequest}$>$ & Entries to create. \\ \hline
\end{tabularx}



\msubsection{struct}{BlacklistCreateRequest}

One system to blacklist.


\noindent\begin{tabularx}{\textwidth}{| >{\raggedright\arraybackslash}p{3.9cm} | >{\raggedright\arraybackslash\small}p{5.3cm} | X |} \hline
\rowcolor{gray!33} Field & Type & Description \\ \hline
system\allowbreak{}Name & \pref{Name} & System to blacklist. \\ \hline
expires\allowbreak{}At & \pref{DateTime} & UTC expiry time; no value means no expiry. \\ \hline
reason & \pref{String} & Reason of the blacklisting. \\ \hline
\end{tabularx}



\msubsection{struct}{BlacklistEntry}

A blacklist entry.


\noindent\begin{tabularx}{\textwidth}{| >{\raggedright\arraybackslash}p{3.9cm} | >{\raggedright\arraybackslash\small}p{5.3cm} | X |} \hline
\rowcolor{gray!33} Field & Type & Description \\ \hline
system\allowbreak{}Name & \pref{Name} & The blacklisted system. \\ \hline
created\allowbreak{}By & \pref{Name} & System that created the entry. \\ \hline
revoked\allowbreak{}By & \pref{String} & System that inactivated the entry, if any. \\ \hline
created\allowbreak{}At & \pref{DateTime} & UTC time of creation. \\ \hline
updated\allowbreak{}At & \pref{DateTime} & UTC time of the last change. \\ \hline
reason & \pref{String} & Reason of the blacklisting. \\ \hline
expires\allowbreak{}At & \pref{DateTime} & Expiry of the entry; no value means permanent. \\ \hline
active & \pref{Boolean} & Whether the entry is in force. \\ \hline
\end{tabularx}



\msubsection{struct}{BlacklistEntryListResponse}

List of blacklist entries.


\noindent\begin{tabularx}{\textwidth}{| >{\raggedright\arraybackslash}p{3.9cm} | >{\raggedright\arraybackslash\small}p{5.3cm} | X |} \hline
\rowcolor{gray!33} Field & Type & Description \\ \hline
entries & \pref{List}$<$\allowbreak{}\mref{BlacklistEntry}$>$ & Matching records of the current page. \\ \hline
count & \pref{Number} & Total number of matching records (all pages). \\ \hline
\end{tabularx}



\msubsection{struct}{BlacklistQueryRequest}

Paged filter for blacklist entries.


\noindent\begin{tabularx}{\textwidth}{| >{\raggedright\arraybackslash}p{3.9cm} | >{\raggedright\arraybackslash\small}p{5.3cm} | X |} \hline
\rowcolor{gray!33} Field & Type & Description \\ \hline
pagination & \mref{Page} & Paging and sorting of the result. \\ \hline
system\allowbreak{}Names & \pref{List}$<$\allowbreak{}\pref{Name}$>$ & System names to match. \\ \hline
mode & \pref{String} & Which entries to return: \texttt{ALL}, \texttt{ACTIVES} or \texttt{INACTIVES}. \\ \hline
issuers & \pref{List}$<$\allowbreak{}\pref{String}$>$ & Systems that created the entries. \\ \hline
revokers & \pref{List}$<$\allowbreak{}\pref{String}$>$ & Systems that inactivated the entries. \\ \hline
reason & \pref{String} & Reason of the blacklisting. \\ \hline
alives\allowbreak{}At & \pref{DateTime} & Only entries in force at this time. \\ \hline
\end{tabularx}



\msubsection{struct}{Page}

Paging and sorting; all fields optional.


\noindent\begin{tabularx}{\textwidth}{| >{\raggedright\arraybackslash}p{3.9cm} | >{\raggedright\arraybackslash\small}p{5.3cm} | X |} \hline
\rowcolor{gray!33} Field & Type & Description \\ \hline
page & \pref{Number} & Zero-based page number. \\ \hline
size & \pref{Number} & Number of records per page (at most \texttt{max.\allowbreak{}page.\allowbreak{}size}, default 1000). \\ \hline
direction & \pref{String} & Sort direction, \texttt{ASC} or \texttt{DESC}. \\ \hline
sort\allowbreak{}Field & \pref{String} & Field to sort by. \\ \hline
\end{tabularx}



\subsection{Primitives}
\label{sec:model:primitives}

Types and structures mentioned throughout this document that are assumed to be available to implementations of this service.
The concrete interpretations of each of these types and structures must be provided by any IDD document claiming to implement this service.

\noindent\begin{tabularx}{\textwidth}{| p{4.6cm} | X |} \hline
\rowcolor{gray!33} Type & Description \\ \hline
\pdef{Boolean} & One out of \texttt{true} or \texttt{false}. \\ \hline
\pdef{DateTime} & Pinpoints a specific moment in time (UTC). \\ \hline
\pdef{List} & An \textit{array} of a known number of items, each having type A. \\ \hline
\pdef{Name} & A string identifier that is intended to be both human and machine-readable. \\ \hline
\pdef{Number} & An integer number. \\ \hline
\pdef{String} & A chain of characters. \\ \hline
\end{tabularx}

\newpage

\bibliographystyle{IEEEtran}
\bibliography{bibliography}

\newpage

\section{Revision History}
\subsection{Amendments}

\noindent\begin{tabularx}{\textwidth}{| p{1cm} | p{3cm} | p{2cm} | X | p{4cm} |} \hline
\rowcolor{gray!33} No. & Date & Version & Subject of Amendments & Author \\ \hline

1 & 2026-10-22 & \arrowversion & Document created from the Arrowhead LaTeX templates for the 5.2.2 Java Spring implementation, including the local extensions of this workspace. & Jerker Delsing \\ \hline
2 & 2026-10-30 & \arrowversion & Checked against the primary Java source tree (Eclipse-Arrowhead-java): services, operations, paths and data models unchanged. & Claude Opus 5.5 (draft) \\ \hline
\end{tabularx}

\subsection{Quality Assurance}

\noindent\begin{tabularx}{\textwidth}{| p{1cm} | p{3cm} | p{2cm} | X |} \hline
\rowcolor{gray!33} No. & Date & Version & Approved by \\ \hline

1 & & \arrowversion  &  \\ \hline

\end{tabularx}

\end{document}
